
Prior work on belief-space planning can be broadly categorized into two classes. Approximate belief-space methods improve scalability by simplifying uncertainty propagation or optimizing surrogate objectives, but generally sacrifice global optimality guarantees \cite{platt2010belief}, \cite{theurkauf2025coop}. In contrast, methods that retain stronger theoretical guarantees typically rely on restrictive assumptions such as full observability, limited coupling between agents, or simplified noise models.

Multi-robot planning methods range from centralized to decoupled. Centralized joint belief-space planning can in principle yield globally optimal solutions but suffers from exponential growth of the state space. Conflict-Based Search (CBS) is a well-known centralized MAPP solver that splits conflicts between agent paths to improve efficiency \cite{sharon2021cbs}, though it operates in a purely discrete setting without uncertainty.  In contrast, decoupled and prioritized planners plan independently for each robot and resolve conflicts or coordinate online, improving scalability but often ignoring dynamical coupling or estimation effects.  Distributed schemes and formation controllers are also used in cooperative localization to maintain structured geometries that reduce drift, but these heuristics generally lack formal guarantees on global optimality.

Grid-based and sampling-based planners provide practical tools for handling large environments under kinematic constraints.  Hexagonal grid discretizations reduce directional bias and provide uniform neighbor spacing, leading to more isotropic path representations \cite{duszak2021hexgrid}.  We adopt a 2D hex grid for the discrete search to better approximate curvature constraints and reduce quantization error. After a discrete solution is obtained, continuous refinement enforces dynamic feasibility. Spline-based trajectory representations, particularly B-splines, are widely used due to their smoothness and local control. Recent works integrate B-splines into sampling-based planners to satisfy curvature limits; for example, variants of RRT* generate tree expansions composed of curvature-constrained B-spline segments \cite{feng2024bsrrt,karaman2011optimal}. These methods embody a hybrid discrete-to-continuous paradigm, combining global search with smooth trajectory optimization.

Cooperative localization and uncertainty propagation are typically handled using distributed estimation frameworks. Kalman filters and their variants propagate agent state covariances, while covariance intersection (CI) enables consistent fusion when cross-correlations are unknown \cite{julier1997cov}. Consensus-based filters allow agents to iteratively exchange estimates with neighbors \cite{olfati2007consensus}.  Extensions often incorporate adaptive or distance-dependent weighting of inter-agent updates to account for communication quality and information relevance.  In particular, spatially decaying Gaussian weights are commonly used to model how sensing and communication efficacy diminish with distance \cite{theurkauf2025coop}. These principles motivate our use of distance-weighted Kalman updates during planning to bound the growth of localization uncertainty under communication constraints.
